\documentclass[10pt,a4paper]{amsart}
\usepackage[margin=19mm]{geometry}
\usepackage{amsmath,amssymb,mathtools}
\usepackage[scr=rsfs,cal=euler]{mathalfa}
\usepackage{xfrac}
\usepackage[unicode,hidelinks,bookmarks=false]{hyperref}
\newcommand{\ie}{i.e.\ }
\newcommand{\cf}{cf.\ }
\newcommand{\resp}{respectively\ }
\newcommand{\Subsection}{\subsection}
\providecommand{\nobreakdash}{\nobreak\hbox{-}}
\setlength{\emergencystretch}{2em}
\multlinegap0pt
\usepackage{bm}
\usepackage{bbm}
\usepackage{dsfont}
\usepackage{xspace}
\usepackage{cancel}
\usepackage{mathtools}
\usepackage{amsthm}
\makeatletter
\@ifundefined{theorem}{\newtheorem{theorem}{Theorem}}{}
\makeatother
\title[Dots and Boxes on Certain Families of Graphs]{Dots and Boxes on Certain Families of Graphs}
\author{Vedant Aryan, Alana Palmer, Alexander Skula, Matthew Woolbert, Joshua Zelinsky}
\date{Source: arXiv:2407.15198v2 (CC BY 4.0)}
\begin{document}
\maketitle

\noindent\emph{Source and licence.} Dots and Boxes on Certain Families of Graphs. Version arXiv:2407.15198v2 (CC BY 4.0), distributed under the Creative Commons Attribution 4.0 International licence, as recorded in the arXiv licence metadata for this exact version and shown on the abstract page https://arxiv.org/abs/2407.15198v2. Full rights notice: licenses/arxiv-2407.15198-CC-BY-4.0.txt.\par\smallskip

\noindent\textit{Editorial scope.} Counted unit: the Theorem whose source label is \texttt{None} in the article (source0.tex lines 674--676, source label \texttt{None}), delivered together with the article's own complete original proof of it (source0.tex lines 678--701, one \texttt{proof} environment of 24 lines). No mathematics is written, repaired or re-derived: statement and proof are the article's own text line for line. Required context retained uncounted: none. Every definition, display and label of those lines is retained unchanged. Omitted from this package and not redistributed: 1--673, 677--677, 702--1167. Nothing inside a retained excerpt is omitted or altered: every retained body is delivered exactly as the article's lines are written, so no omission receipt of any kind accompanies this package. Citations: the retained excerpts carry no citation command at all, so no bibliography entry of the article is transcribed and no reference is redistributed. Reference adaptation: none was needed and none was made; every reference inside the delivered unit resolves inside it, so no unresolved reference or citation is produced. Printed slips kept literal: no printed word, symbol, hypothesis or number of the counted statement or of its proof was altered. Layout and class: the article's own class is \texttt{article} with packages including caption, graphicx, amsthm, amsfonts, amssymb, amsmath, epsf, verbatim, cite, tikz, url, placeins, geometry; the wrapper is the lane's pinned \texttt{amsart} editorial wrapper at 10pt on A4, which restates the theorem-like environments the retained text needs and adds the article's own macro definitions that the retained text needs (none), with extra wrapper packages (bm, bbm, dsfont, xspace, cancel, mathtools). The delivered numbering is the unit's own. Assets: no retained span contains a figure environment, a graphics-inclusion command, a PDF-page inclusion, a table or a TikZ picture; no image file is copied into this package, the package declares no asset dependency and no image file is redistributed with it. Licence: version arXiv:2407.15198v2 of this article is distributed under the Creative Commons Attribution 4.0 International licence, as recorded in the arXiv licence metadata for this exact version and shown on the abstract page https://arxiv.org/abs/2407.15198v2. A full rights notice is delivered at licenses/arxiv-2407.15198-CC-BY-4.0.txt. Characters: every retained excerpt is pure ASCII, so the delivered engine, XeLaTeX with its default Latin Modern text font, reports no missing character.
\begin{theorem}
For all complete bipartite graphs $K(2, b)$ with $b \geq 2$, Player 1 wins when $b$ is odd and Player 2 wins when $b$ is even. Specifically, if $b = 2k$ is even, then the score is $(k-1, k+3)$ so Player 2 wins by 4 points. If $b = 2k+1$ is odd, then the score is $(k+3, k)$, so Player 1 wins by 3 points.
\end{theorem}

\begin{proof}
We will prove this theorem by induction on $b$.

\textbf{Base cases:}
\begin{itemize}
    \item $K(2, 2)$: This graph is isomorphic to a cycle $C_4$. Player 2 wins with a score of $(0, 4)$, which matches our formula with $k = 1$: $(k-1, k+3) = (0, 4)$.
    
    \item $K(2, 3)$: After Player 1's first move, one vertex in the second partite set will have degree 1. Player 2 must take the edge connecting to this vertex, gaining the vertex. The remaining graph is $K(2, 2)$, which is a Player 2 win with score $(0, 4)$. Thus, the overall score for $K(2, 3)$ is $(3, 1) + (0, 4) = (3, 1)$, which matches our formula with $k = 1$: $(k+3, k) = (4, 1)$.
\end{itemize}

\textbf{Inductive step:} Assume that for some $n \geq 3$, the theorem holds for all values of $b$ such that $2 \leq b < n$. We want to show it holds for $b = n$.

Consider the game on $K(2, n)$. When Player 1 makes the first move, they must take an edge without gaining any vertex (since all vertices in the second partite set have degree 2 initially). This leaves one vertex in the second partite set with degree 1. Player 2 then takes the edge connecting to this vertex, gaining the vertex.

The remaining graph is $K(2, n-1)$. By the induction hypothesis:

\begin{itemize}
    \item If $n-1 = 2k$ is even, then the score on $K(2, n-1)$ is $(k-1, k+3)$. Player 2 has already gained one vertex, so the overall score is $(k-1, k+3+1) = (k-1, k+4)$. Substituting $n = 2k+1$, we get a score of $(k-1, k+4) = ((n-1)/2 - 1, (n-1)/2 + 4) = ((n-3)/2, (n+7)/2)$. For $n = 2k+1$, this becomes $(k-1, k+4) = (k-1, (2k+1+7)/2) = (k-1, k+4)$, which matches our formula for odd $b$.
    
    \item If $n-1 = 2k+1$ is odd, then the score on $K(2, n-1)$ is $(k+3, k)$. Player 2 has already gained one vertex, so the overall score is $(k+3, k+1)$. Substituting $n = 2k+2$, we get a score of $(k+3, k+1) = (n/2 + 2, n/2 - 1)$. For $n = 2k+2$, this becomes $(k+3, k+1) = ((2k+2)/2 + 2, (2k+2)/2 - 1) = (k+3, k+1) = ((n-2)/2 + 3, (n-2)/2 - 1 + 2) = ((n+4)/2, (n+2)/2 - 1)$, which matches our formula for even $b$.
\end{itemize}

Therefore, by the principle of mathematical induction, the theorem holds for all $b \geq 2$.
\end{proof}

\end{document}
